% generated by sdr.manuscript; do not edit
{\footnotesize
\begin{longtable}{>{\raggedright\arraybackslash}p{\dimexpr0.0745\linewidth-2\tabcolsep\relax}>{\raggedright\arraybackslash}p{\dimexpr0.1383\linewidth-2\tabcolsep\relax}>{\raggedright\arraybackslash}p{\dimexpr0.4362\linewidth-2\tabcolsep\relax}>{\raggedright\arraybackslash}p{\dimexpr0.3511\linewidth-2\tabcolsep\relax}}
\caption{Eligibility criteria (PROTOCOL.md section 4).}\label{tab:eligibility}\\
\toprule
\textbf{Code} & \textbf{Criterion} & \textbf{Decision rule} & \textbf{Rationale} \\ \midrule
\endfirsthead
\multicolumn{4}{l}{\emph{\tablename~\thetable{} (continued)}}\\
\toprule
\textbf{Code} & \textbf{Criterion} & \textbf{Decision rule} & \textbf{Rationale} \\ \midrule
\endhead
\midrule\multicolumn{4}{r}{\emph{continued}}\\
\endfoot
\bottomrule
\endlastfoot
I1 & Human spine imaging & The description says the images show the human spine or a named vertebral region (cervical, thoracic, lumbar, sacrum, whole spine), or the spine is a labelled structure in a body-region dataset (a whole-body CT dataset with vertebra masks passes). Chest, abdomen or pelvis datasets with no spine annotation and no spine-specific description are \textbf{not} included merely because vertebrae are in the field of view. & Vertebrae sit in the field of view of most chest, abdomen and pelvis scans; without a spine claim or a spine label such a dataset cannot be found by a spine search or charted for annotation origin, and the count of "spine datasets" would become a count of body CT. \\
I2 & Modality & At least one of: CT, CBCT, plain radiograph (CR/\allowbreak{}DR/\allowbreak{}film), fluoroscopy or C-arm frames, DXA vertebral fracture assessment images, EOS biplanar radiographs. Multi-modal datasets are included when at least one eligible modality has at least 10 subjects; ineligible modalities are recorded but not counted. & These are the modalities of the intended use (level identification from projection images and from CT). CBCT is included because intra-operative CBCT is the CT most often paired with fluoroscopy. MRI has its own label conventions and literature and is not part of the question. \\
I3 & Size & I3 holds when a stated \texttt{n\_\allowbreak{}subjects} $\geq$ 10, \textbf{or} when \texttt{n\_\allowbreak{}subjects} is unknown and a stated \texttt{n\_\allowbreak{}images\_\allowbreak{}or\_\allowbreak{}volumes} $\geq$ 10 (charted \texttt{n\_\allowbreak{}subjects\_\allowbreak{}basis =\allowbreak{} unknown}). Figure-derived collections use source articles as the proxy when the subject count is not stated. & Ten is the smallest set for which a per-subject variant count (RQ3) is a count rather than a case report; the image fallback keeps datasets described only in scans or studies, which repository records often are. \\
I4 & Described & At least one of: journal article, preprint, challenge web page, data-repository record (Zenodo, Figshare, Dryad, OSF, Mendeley Data, PhysioNet, TCIA, Kaggle, Hugging Face, GitHub, institutional repository). A bare download link with no text describing the content fails I4. Authors'-own datasets are no exception: FluoroCap, or any other set an author maintains, enters only if a public describing record exists on the run date; otherwise it is reported as "known item, excluded I4" in the other-methods arm of the flow (section 7.\allowbreak{}7). & This review charts only public text; a dataset a reader cannot find a description of cannot be charted from it. \\
I5 & Date & The earliest record that describes the \emph{image release} (article, preprint, challenge page or repository record) is dated 2000-01-01 to 2026-12-31. \texttt{year} is that record's year; the years in which the images were acquired go in \texttt{cohort\_\allowbreak{}year\_\allowbreak{}range} and do not decide I5 (NHANES II films, acquired 1976--80 and released as digital images by NLM in the 2000s, pass I5 on the release record). & 2000 is the first year in which spine image collections were distributed digitally rather than as film libraries; earlier cohorts enter through their digital release. \\
I6 & Annotation status & Any: with annotations or without. & Unannotated datasets are what the project's own pipeline labels; excluding them would answer RQ1 wrongly. \\
I7 & Figure-derived collections & Collections assembled from published figures (PMC, Open-i, MultiCaRe, Open-PMC-18M, Wikimedia Commons, journal APIs) are eligible when I1--I5 hold and are flagged \texttt{figure\_\allowbreak{}derived =\allowbreak{} true}. & The only public source of fluoroscopy at scale, and what FluoroCap was built from. \\
I8 & Derived datasets & Datasets built by re-annotating, re-labelling, subsetting or converting another public dataset are eligible, flagged \texttt{derived =\allowbreak{} true}, with the parent recorded in \texttt{parent\_\allowbreak{}ids}. & Re-annotations are where level names and variant labels are most often added to an older release. \\
E1 & MRI only & No eligible modality is present (MRI-only, ultrasound-only, nuclear-medicine-only). & Out of scope by I2. \\
E2 & Non-human & Animal, veterinary or comparative-anatomy imaging. A dataset mixing human and animal subjects is included with the human count only if that count is stated; otherwise excluded. & Population is the human spine. \\
E3 & Phantom or cadaver only & All images are of phantoms, cadavers, ex-vivo specimens or printed models. A dataset mixing living subjects and cadavers is included with the living count. & Variant prevalence and level naming are questions about living anatomy in clinical images. \\
E4 & Purely synthetic & Images generated (DRR, GAN, diffusion, simulation) with no public parent CT or radiograph dataset. Synthetic datasets rendered from a public parent are included under I8 as \texttt{derived}. & A synthetic set with no public parent has no subjects to count and no annotation origin to chart. \\
E5 & Too small & A \emph{stated} subject count below 10. A record with an unknown subject count and at least 10 images or volumes does not fail E5 (see I3). & The complement of I3. \\
E6 & Not subject-level images & The release contains only extracted features, measurements, embeddings, landmarks without pixels, reports or metadata, without the images. & The review is of image datasets. \\
E7 & Duplicate record & \textbf{Not a screening code.\allowbreak{}} A second, third \ldots{} record about an already included dataset is not excluded at screening: it passes to record-to-dataset mapping (section 7.\allowbreak{}6) with a role, so that it stays in \texttt{record\_\allowbreak{}ids}, in backward chasing and in the source-overlap diagnostic (SAP A9). Versions of one dataset collapse to one dataset with the versions listed; a derived dataset is not a duplicate of its parent. Identical records (same canonical identifier, section 7.\allowbreak{}1) are removed at de-duplication and counted there. & Dropping later records at screening silently removes the evidence the flow and the overlap diagnostic count. \\
E8 & Not a dataset & A method paper that uses a private dataset without describing it as available; a review, editorial, protocol; a model or software release without images. A method paper that describes and releases its data passes. A method paper that merely \emph{uses} an already public dataset and releases no new images or labels is coded \texttt{use-\allowbreak{}only} (section 7.\allowbreak{}3): not an included record, counted in that dataset's \texttt{n\_\allowbreak{}use\_\allowbreak{}records}. Reviews are retained until charting for backward chasing (section 6.\allowbreak{}4), then excluded under E8. & Use-only papers (VerSe alone has hundreds) would otherwise inflate the included-record count without adding a dataset. \\
\end{longtable}}
